% GENERATED FILE. Edit canonical lessons, then run npm run generate:books.
% canonical-source-hash: fnv1a64:447e2e1cbb8c83a9
% canonical-lessons: ES-C56-si-tuviera

\chapter{Si Tuviera --- Supposing What Is Not So}
\label{ch:si-tuviera}
\hlchaptermodality{208}

\begin{chapteropening}
\textbf{By the end of this chapter:} \emph{I can say what would happen if something untrue were true.}

I can say what would happen if something untrue were true.
\end{chapteropening}

\section[si tuviera]{si tuviera}

\label{lesson:ES-C56-si-tuviera}



\begin{quote}
Say the real one. (\emph{Si tengo tiempo, hablo español}.) Now the supposing form of \emph{tener}. (\emph{Tuviera}.)

Put them together and you can talk about a day you are not having.
\end{quote}

\begin{grammarlens}[title={Grammar Lens: the world that is not the case}]
\begin{quote}
\textbf{si + supposing form, then the \emph{-ría} form.}
\end{quote}

\noindent
\begin{tabularx}{\linewidth}{@{}>{\raggedright\arraybackslash}X>{\raggedright\arraybackslash}X@{}}
\toprule
\textbf{} & \textbf{} \\
\midrule
\emph{Si tengo tiempo, hablo español.} & and I might well have it \\
\emph{Si tuviera tiempo, hablaría español.} & \textbf{but I do not} \\
\bottomrule
\end{tabularx}

The second sentence says two things at once: what would happen, \textbf{and} that it is not happening. That is carried entirely by the verb forms — no extra word is needed to say \emph{but I don't}.

Both halves change together. \emph{Si} takes the supposing form; the result half takes the \emph{-ría} form you learned for \emph{would}. Neither half can be left in the present.
\end{grammarlens}

\begin{culture}
\textbf{And \emph{si} never takes the present subjunctive.} Not \emph{si tenga} — that combination simply does not occur in Spanish. The supposing form after \emph{si} is always the \emph{-ra} one.

It is worth knowing why the pairing feels lopsided. The \emph{si} half is not asserting anything, so it uses the form Spanish reserves for the unasserted — and the result half \textbf{is} making a claim, just a claim about a world that does not exist, which is what \emph{-ría} is for. Two different jobs, two different forms, in one sentence.
\end{culture}

\subsection*{Your turn}
Say these aloud:

\begin{itemize}
\raggedright
  \item \emph{Si tuviera tiempo, hablaría español}
  \item the real version, then the unreal one, back to back
\end{itemize}

\emph{Twice through:}

\begin{tcolorbox}[breakable,colback=teal!4,colframe=teal!35!black,title={Before you move on}]
Which form follows \emph{si} here? (The \textbf{supposing} form — \emph{tuviera}.) Which follows the comma? (The \emph{\textbf{-ría}} form.)

And what does the sentence say besides its words? (That the thing \textbf{is not true}.)
\end{tcolorbox}
